% ch05  参考文献（原文 2029–2276 行）
% 由脚本机械转换：条目标题保留英文，序号与原文一致。

\clearpage
\phantomsection
\addcontentsline{toc}{section}{参考文献}

\begin{thebibliography}{99}

\bibitem{ARH}
E. Arasteh Rad, U. Hartl,
Uniformizing the moduli stacks of global G-shtukas.
Int. Math. Res. Not. IMRN 2021, no. 21, 16121--16192.

\bibitem{And}
G.Anderson, 
$t$-motives.
Duke Math. J. 53 (1986), no. 2, 457--502.

\bibitem{Br}
P. Breutmann, 
Functoriality of Moduli Spaces of Global $G$-Shtukas. 
arXiv:1902.10602.

\bibitem{DL}
P. Deligne, G. Lusztig,
Representations of reductive groups over finite fields. 
Ann. of Math. (2) 103 (1976), no. 1, 103--161.

\bibitem{Dr-Ell}
V. Drinfeld,
Elliptic modules. Mat. Sb. (N.S.) 94(136) (1974), 594--627, 656. Elliptic modules II. Mat. Sb. (N.S.) 102(144) (1977), no. 2, 182--194, 325.

\bibitem{Dr-FBun}
V. Drinfeld,
Moduli varieties of $F$-sheaves. 
Funktsional. Anal. i Prilozhen. 21 (1987), no. 2, 23--41. 
 
\bibitem{Dr-Pet}
V. Drinfeld, 
Proof of the Petersson conjecture for $\GL(2)$ over a global field of characteristic $p$. 
Funktsional. Anal. i Prilozhen. 22 (1988), no. 1, 34--54, 96; translation in Funct. Anal. Appl. 22 (1988), no. 1, 28--43.

\bibitem{Dr-comp}
V.Drinfeld, 
Cohomology of compactified moduli varieties of $F$-sheaves of rank 2.  
Zap. Nauchn. Sem. Leningrad. Otdel. Mat. Inst. Steklov. (LOMI) 162 (1987), Avtomorfn. Funkts. i Teor. Chisel. III, 107--158, 189; translation in J. Soviet Math. 46 (1989), no. 2, 1789--1821.

\bibitem{FYZ1}
T. Feng, Z. Yun, W. Zhang,
Higher Siegel--Weil formula for unitary groups: the non-singular terms.
Invent. Math. 235 (2024), no. 2, 569--668.

\bibitem{FYZ2}
T. Feng, Z. Yun, W. Zhang,
Higher theta series for unitary groups over function fields.
arXiv:2110.07001. To appear in Ann. Sci. E.N.S.

\bibitem{FYZ3}
T. Feng, Z. Yun, W. Zhang,
Modularity of higher theta series I: cohomology of the generic fiber.
arXiv: 2308.10979.

\bibitem{GGP}
W-T. Gan, B.Gross, D.Prasad,
Symplectic local root numbers, central critical $L$ values, and restriction problems in the representation theory of classical groups. 
Sur les conjectures de Gross et Prasad. I. Ast\'erisque No. 346 (2012), 1--109.

\bibitem{GL}
A. Genestier, V. Lafforgue,
Chtoucas restreints pour les groupes r\'eductifs et param\'etrisation de Langlands locale.
arXiv:1709.00978.

\bibitem{Goss}
D. Goss, 
Basic structures of function field arithmetic.
Ergebnisse der Mathematik und ihrer Grenzgebiete (3), 35. Springer-Verlag, Berlin, 1996. xiv+422 pp.

\bibitem{Jan}
U. Jannsen,
Motives, numerical equivalence, and semi-simplicity, 
Invent. Math. 107, 1992, 447--452.

\bibitem{Kot}
R. Kottwitz,
$B(G)$ for all local and global fields.
arXiv:1401.5728.

\bibitem{Kud}
S.Kudla,
Central derivatives of Eisenstein series and height pairings, 
Ann. of Math. (2) 146 (1997), no. 3, 545--646.

\bibitem{KR1}
S. Kudla, M.Rapoport,
Special cycles on unitary Shimura varieties I: Unramified local theory, 
Invent. Math. 184 (2011), no. 3, 629--682.

\bibitem{KR2}
S. Kudla, M. Rapoport, 
Special cycles on unitary Shimura varieties II: Global theory, 
J. Reine Angew. Math. 697 (2014), 91--157.

\bibitem{Laff-comp}
L. Lafforgue,
Une compactification des champs classifiant les chtoucas de Drinfeld. 
J. Amer. Math. Soc. 11 (1998), no. 4, 1001--1036.

\bibitem{LLaff}
L. Lafforgue,
Chtoucas de Drinfeld et correspondance de Langlands. 
Invent. Math. 147 (2002), no. 1, 1--241.

\bibitem{VLaff}
V. Lafforgue,
Chtoucas pour les groupes r\'eductifs et param\'etrisation de Langlands globale.
J. Amer. Math. Soc. 31 (2018), no. 3, 719--891.

 
\bibitem{LR}
R. Langlands, M.Rapoport,
Shimuravariet\"aten und Gerben. 
J. Reine Angew. Math. 378 (1987), 113--220.

\bibitem{LRS}
G. Laumon, M. Rapoport, U. Stuhler,
$\cD$-elliptic sheaves and the Langlands correspondence.
Invent. Math. 113 (1993), no. 2, 217--338.

\bibitem{Laumon}
G. Laumon,
Drinfeld shtukas. 
In {\em Vector bundles on curves-- new directions} (Cetraro, 1995), 50--109, Lecture Notes in Math., 1649, Fond. CIME/CIME Found. Subser., Springer, Berlin, 1997.

\bibitem{LZ}
C. Li, W. Zhang,
Kudla-Rapoport cycles and derivatives of local densities. 
J. Amer. Math. Soc. 35 (2022), no. 3, 705--797. 

\bibitem{Milne}
J. Milne, 
Motives over finite fields.
In {\em Motives} (Seattle, WA, 1991), 401--459, 
Proc. Sympos. Pure Math., 55, Part 1, Amer. Math. Soc., Providence, RI, 1994.

\bibitem{MV}
I. Mirkovic, K.Vilonen,
Geometric Langlands duality and representations of algebraic groups over commutative rings. 
Ann. of Math. (2) 166 (2007), no. 1, 95--143.

\bibitem{Mum}
D. Mumford,
An algebro-geometric construction of commuting operators and of solutions to the Toda lattice equation, Korteweg deVries equation and related nonlinear equation. 
In {\em Proceedings of the International Symposium on Algebraic Geometry} (Kyoto Univ., Kyoto, 1977), pp. 115--153, Kinokuniya Book Store, Tokyo, 1978.

\bibitem{NDT}
T. Ng\^o Dac,
Compactification des champs de chtoucas et th\'eorie g\'eom\'etrique des invariants.
Ast\'erisque No. 313 (2007), 124 pp.

\bibitem{NDT-Sht}
T. Ng\^o Dac,
Introduction to the stacks of shtukas. 
Algebraic cycles, sheaves, shtukas, and moduli, 217--236, 
Trends Math., Birkh\"auser, Basel, 2008.
 
\bibitem{PR}
G. Pappas, M. Rapoport, 
Twisted loop groups and their affine flag varieties. With an appendix by T. Haines and Rapoport. Adv. Math. 219 (2008), no. 1, 118--198. 

\bibitem{Ra}
M. Rapoport,
A guide to the reduction modulo $p$ of Shimura varieties, 
in Automorphic Forms. I, 298, 2005, pp. 271--318.

\bibitem{S-Berkeley}
P. Scholze, J. Weinstein, 
Berkeley lectures on $p$-adic geometry. 
Annals of Mathematics Studies, 207. 
Princeton University Press, Princeton, NJ, 2020. x+250 pp.

\bibitem{SchICM}
P. Scholze, 
$p$-adic geometry. 
In {\em Proceedings of the International Congress of Mathematicians--Rio de Janeiro 2018}. Vol. I. Plenary lectures, 899--933, World Sci. Publ., Hackensack, NJ, 2018.

\bibitem{Var}
Y.Varshavsky, 
Moduli spaces of principal $F$-bundles. 
Selecta Math. (N.S.) 10 (2004), no. 1, 131--166. 

\bibitem{X}
C.Xue,
Finiteness of cohomology groups of stacks of shtukas as modules over Hecke algebras, and applications.
\'Epijournal G\'eom. Alg\'ebrique 4 (2020), Art. 6, 42 pp.

\bibitem{Xue}
C. Xue,
Lectures in the IHES summer school.

\bibitem{YZ1}
Z. Yun, W. Zhang, 
Shtukas and the Taylor expansion of $L$-functions. 
Ann. of Math. (2) 186 (2017), no. 3, 767--911.

\bibitem{YZ2}
Z. Yun, W. Zhang, 
Shtukas and the Taylor expansion of $L$-functions (II). 
Ann. of Math. (2) 189 (2019), no. 2, 393--526. 

\bibitem{Y-ICM}
Z. Yun, 
Hitchin type moduli stacks in automorphic representation theory. 
Proceedings of the International Congress of Mathematicians -- Rio de Janeiro 2018. Vol. II. Invited lectures, 1447--1476, World Sci. Publ.,
Hackensack, NJ, 2018.

\bibitem{Y-cycle}
Z. Yun,
Special cycles for Shtukas are closed.
Pure Appl. Math. Q. 18 (2022), no. 5, 2203--2220.

\bibitem{Zhu-Satake}
X. Zhu, 
An introduction to affine Grassmannians and the geometric Satake equivalence.
In {\em Geometry of Moduli Spaces and Representation Theory}, IAS/Park City Mathematical Series vol 24.

\bibitem{Zhu}
X. Zhu,
Coherent sheaves on the stack of Langlands parameters.
arXiv:2008.02998.

\end{thebibliography}
